\chapter{Notes diverses}

Ce chapitre contient des \'el\'ements de textes partiellement redondant ou
non ais\'ement int\'egrales au reste du m\'emoire\footnote{%
  En autre par manque de temps.} %
et des notes souvent sous forme de questions brute. %
Nous pensons que la pr\'esence des premi\`eres sections au moins,
permettra d'\'eclairer notre d\'emarche pour certaines questions soulev\'ees
dans le reste de ce travail.

\section{lambda}

Il faut pond\'erer le calcul de lambda sur une r\`egle par l'application
de lambda sur l'index de chacune des entr\'ees. %

\section{anneal et life}

L'\'etude de s\'eries temporelles associ\'e \`a anneal. %
Mesure de la population et de la taille des fronti\`eres. %
La taille des fronti\`eres est obtenu par une diff\'erence spatiale. %
Les s\'eries obtenu ne pr\'esentent pas de caract\`ere fractal sur une
grande \'echelle\footnote{%
  La question de la nature du \aquote{bruit} \`a
  petite \'echelle reste pos\'e~(comparer avec les s\'eries tir\'e de life)} %
elles pr\'esentent par contre une invariance au changement d'\'echelle
d'espace, \ie Les s\'erie tendent \`a \^etre identique \`a l'\'echelle du
transient, a une dilatation du temps pr\`es, lorsque la taille de
l'espace tends vers l'infini.

Life filtr\'e~(la seconde diff\'erence spatiale) est suppos\'e fournir des
s\'eries dot\'e d'une dimension fractale sur une grande \'echelle, mais pas
d'invariance au changement d'\'echelle d'espace\footnote{%
  Sauf pour la dimension fractale \'evidemment}.

la quantit\'e de bruit n\'ecessaire pour faire tendre le transient vers
l'infini et suppos\'e petite et d\'ecroissante fonction de la taille de
l'univers. %
Pour anneal le comportement inverse est attendu.

\section{le contr\^ole lin\'eaire des r\`egle}

contr\^ole lin\'eaire de la moyenne?
\begin{itemize}
\item param\'etrage global de comportement
\item formulation diff\'erentielle
\item 1 plan m\'emoire par param\`etre
  et 1 plan m\'emoire pour les $n$ variables\footnote{%
    Structure du mixage spatial?}, ou $n$ plan m\'emoire pour les $n$
  variables et \'equation bool\'eenne donne la piste d'une technique de
  transformation d'\'equation diff\'erentielle continue en \'equation
  bool\'eenne combin\'e.
\end{itemize}
La dotation en espace\footnote{%
  Dimension plus voisinage.}, %
d'une \'equation diff\'erentielle, associ\'e \`a sa bool\'ea\-ni\-sa\-tion permet de
garantir sa calculabilit\'e.

\section{Espace, temps, lois}

L'espace est constitu\'e de cellules \`a \'etat finit dot\'e d'une dimension
et d'un r\'eseau d'interconnexion. %
Il correspond \`a la ressource m\'emoire mesur\'e en \'etendu. %
%%
Le temps est un objet plus difficile \`a appr\'ehender, il correspond \`a
la ressource commutation mesur\'e en d\'ebit, mais se ram\`ene par
plongement \`a une manipulation d'espace, et correspond alors lui aussi
\`a une ressource m\'emoire. %
Il se caract\'erise \'egalement par les op\'erateurs appliqu\'e sur un
espace-temps cellulaire, la r\'eduction op\'er\'e par ceux ci \'etant
g\'en\'eralement une projection perpendiculaire \`a l'axe temporel.
%%
Les lois, ou r\`egles ou fonctions de transitions d\'efinissent des
mapping locaux de de l'espace sur lui m\^eme. %
Le nombre d'\'etats cellulaires et le nombre d'\'el\'ements du voisinage
cellulaire d\'etermine la taille de l'espace de r\`egles associ\'e et le
co\^ut de la transformation d'une r\`egles exprim\'e comme une s\'equences
d'instructions en table de transition.

\section{les AC de dimensions~1}

Nous avons d\'elib\'er\'ement ignor\'e dans ce travail les AC de dimensions 1,
qui constitue le support de la majorit\'e des recherches sur la
structure de l'espace de r\`egles, pour nous int\'eresser aux AC de
dimension 2. %
Ce choix correspond \`a une volont\'e de chercher les limites pratique du
calcul cellulaires sur des architecture mat\'erielles
standard\footnote{%
  Nous entendons par l\`a des r\'eseaux de stations de travail.} %.
Nous pensons que la domination de AC de dimension 1 dans les travaux
de recherche r\'ecent sur la structure de l'espace des r\`egles est du en
partie\footnote{%
  La raison primordiale \'etant a priori la volont\'e de chercher les
  mod\`eles minimum g\'en\'erant les comportement les plus riches possibles,
  et la facilit\'e d'analyse.} %
\`a la relative difficult\'e de r\'ealiser des exp\'erience efficace
sur des espaces de configurations de dimension sup\'erieur \`a~1.

%%\input{tthy-notes-futur}

%%
%% Local Variables:
%% mode: auto-fill
%% iso-accents-minor-mode: nil
%% TeX-master: "top"
%% TeX-command-default: "mytex"
%% Local IspellParsing: "latex-mode ~latin1"
%% Local IspellPersDict: "~/.ispell_lisp"
%% LocalWords: "lookup" lambda anneal fractal fractale param\'etrage bool\'eenne PE
%% End:
% LocalWords:  d'\'equation calculabilit\'e SIMD s\'equenceurs d'AC s\'equencement in
% LocalWords:  sous-AC espace-temps mapping algo debug portabilit\'e
